\documentclass[11pt]{article}

\usepackage[T1]{fontenc}
\usepackage{lmodern}
\usepackage[margin=0.8in]{geometry}
\usepackage{amsmath,amssymb}
\usepackage{booktabs}
\usepackage{array}
\usepackage{xcolor}
\usepackage{titlesec}
\usepackage{lastpage}
\usepackage{fancyhdr}

\definecolor{accent}{HTML}{0F6CBF}
\definecolor{band}{HTML}{6B5BD0}

\titleformat{\section}{\normalfont\large\bfseries\color{accent}}{\thesection.}{0.5em}{}
\titlespacing*{\section}{0pt}{0.4em}{0.2em}

\pagestyle{fancy}
\fancyhf{}
\renewcommand{\headrulewidth}{0pt}
\fancyfoot[C]{\footnotesize\color{gray} Intro Physics II \textbullet\ Circuits \textbullet\ page \thepage\ of \pageref{LastPage}}

\setlength{\parindent}{0pt}
\raggedbottom
\setlength{\extrarowheight}{2pt}
\renewcommand{\arraystretch}{1.05}
\setlength{\aboverulesep}{1pt}\setlength{\belowrulesep}{1pt}

% two-column equation table: left = description, right = formula
\newcolumntype{L}{>{\raggedright\arraybackslash}p{0.40\linewidth}}
\newcolumntype{F}{>{\raggedright\arraybackslash}p{0.555\linewidth}}
\newcommand{\eqtab}[1]{%
  \par\vspace{0.15em}
  \noindent
  \begin{tabular}{@{}L F@{}}
  \toprule
  #1
  \bottomrule
  \end{tabular}
  \par\vspace{0.1em}}
\newcommand{\R}[2]{#1 & $\displaystyle #2$ \\}
% sub-band label spanning both columns
\newcommand{\band}[1]{%
  \addlinespace[1.5pt]%
  \multicolumn{2}{@{}l@{}}{\rule{0pt}{1.0em}\footnotesize\bfseries\color{band}\textsc{#1}}\\[2pt]}

\begin{document}

\begin{center}
  {\LARGE\bfseries Equation Sheet: Circuits}\\[2pt]
  {\large Intro Physics II \quad\textbullet\quad Lectures 14--17 \quad\textbullet\quad Knight ch.\ 26--28}
\end{center}
\vspace{0.25em}

\hrule
\vspace{0.4em}
{\small \textbf{\textsc{\textcolor{band}{fundamental}}} = know cold (your starting points); \textbf{\textsc{\textcolor{band}{derived}}} = a line or two from the fundamentals.}
\vspace{0.25em}
\hrule
\vspace{0.25em}
{\small\bfseries\color{band}\textsc{Symbols and constants}}\\[2pt]
\begingroup
\footnotesize
\setlength{\extrarowheight}{0pt}\renewcommand{\arraystretch}{1.05}
\begin{tabular}{@{}r@{\ \ }l@{\hspace{1.8em}}r@{\ \ }l@{\hspace{1.8em}}r@{\ \ }l@{}}
$I$ & current (A $=$ C/s)        & $\mathcal{E}$ & emf (V)                       & $P$ & power (W) \\
$V,\Delta V$ & potential difference (V) & $r$    & internal resistance ($\Omega$) & $Q$ & charge (C) \\
$R$ & resistance ($\Omega$)      & $C$          & capacitance (F)               & $\tau$ & time constant (s) \\
$\rho$ & resistivity ($\Omega\cdot$m) & $L,A$ & wire length / cross-section area & $e$ & $1.60\times10^{-19}$ C \\
\end{tabular}
\endgroup
\vspace{0.1em}
\hrule

%========================================================
\section{Current, resistance \& Ohm's law \quad {\normalfont\small (Knight 26.1--26.5, 27.1--27.4)}}
\eqtab{
  \band{fundamental}
  \R{Current --- rate of charge flow}{I = \dfrac{dQ}{dt} \quad\xrightarrow{\ \text{steady}\ }\quad I = \dfrac{Q}{\Delta t}}
  \R{Ohm's law\newline\footnotesize(for an ohmic resistor; also $\Delta V = IR$ across any $R$)}{I = \dfrac{\Delta V}{R}}
  \R{Resistor power\newline\footnotesize(rate electrical energy $\to$ heat/light)}{P = I\,\Delta V = I^{2}R = \dfrac{(\Delta V)^{2}}{R}}
  \band{derived}
  \R{Resistance of a uniform wire}{R = \dfrac{\rho L}{A}}
  \R{Conservation of charge at a wire}{I \text{ is the same all along a series path}}
}

%========================================================
\section{EMF \& real batteries \quad {\normalfont\small (Knight 28.2--28.3)}}
\eqtab{
  \band{fundamental}
  \R{EMF --- energy supplied per unit charge}{\mathcal{E} = \dfrac{dW}{dq} \quad (\text{units of volts})}
  \band{derived}
  \R{Terminal voltage of a real battery\newline\footnotesize(internal resistance $r$; equals $\mathcal{E}$ only when $I=0$)}{\Delta V_{\text{term}} = \mathcal{E} - I r}
  \R{Current from a single-loop source}{I = \dfrac{\mathcal{E}}{R + r}}
}

%========================================================
\section{Series \& parallel combinations \quad {\normalfont\small (Knight 28.4)}}
\eqtab{
  \band{fundamental}
  \R{Resistors in series\newline\footnotesize(same current; voltages add; $R_{\text{eq}}$ is bigger than any one)}{R_{\text{eq}} = R_1 + R_2 + \cdots}
  \R{Resistors in parallel\newline\footnotesize(same voltage; currents add; $R_{\text{eq}}$ is smaller than any one)}{\dfrac{1}{R_{\text{eq}}} = \dfrac{1}{R_1} + \dfrac{1}{R_2} + \cdots}
  \band{derived}
  \R{Two resistors in parallel}{R_{\text{eq}} = \dfrac{R_1 R_2}{R_1 + R_2}}
  \R{Capacitors\newline\footnotesize(combine the opposite way to resistors)}{C_{\text{par}} = C_1 + C_2 + \cdots,\qquad \dfrac{1}{C_{\text{ser}}} = \textstyle\sum \dfrac{1}{C_i}}
  \R{Voltage divider (series $R_1,R_2$)}{\Delta V_1 = \Delta V_{\text{tot}}\,\dfrac{R_1}{R_1 + R_2}}
}

%========================================================
\section{Kirchhoff's rules \quad {\normalfont\small (Knight 28.1, 28.5)}}
\eqtab{
  \band{fundamental}
  \R{Junction rule (charge conservation)}{\textstyle\sum I_{\text{in}} = \sum I_{\text{out}}}
  \R{Loop rule (energy conservation)\newline\footnotesize($+$ across a source from $-$ to $+$; $-IR$ across a resistor in the direction of $I$)}{\textstyle\sum_{\text{loop}} \Delta V = 0}
  \band{derived}
  \R{Ideal meters}{\text{ammeter: } R = 0 \text{ (in series)};\ \ \text{voltmeter: } R = \infty \text{ (in parallel)}}
}

%========================================================
\section{RC circuits \quad {\normalfont\small (Knight 28.6)}}
\eqtab{
  \band{fundamental}
  \R{Capacitor relation}{Q = C\,\Delta V_C}
  \R{Time constant}{\tau = RC}
  \band{derived}
  \R{Discharging (starts at $Q_0$, $I_0$)\newline\footnotesize(both fall to $37\%$ after one $\tau$)}{Q(t) = Q_0\,e^{-t/\tau}\,,\qquad I(t) = I_0\,e^{-t/\tau}}
  \R{Charging through a battery $\mathcal{E}$\newline\footnotesize(reaches $63\%$ of full charge after one $\tau$; $I$ decays)}{Q(t) = C\mathcal{E}\left(1 - e^{-t/\tau}\right)}
  \R{Current is largest at $t=0$}{I_0 = \dfrac{\mathcal{E}}{R} \quad(\text{charging});\qquad I_0 = \dfrac{\Delta V_0}{R} \quad(\text{discharging})}
}

\vfill
{\footnotesize\color{gray}\raggedright
Companion simulations at \texttt{hoppese.github.io/Intro-physics-sims/}\,:
\texttt{circuit-builder}, \texttt{circuit-lab}, \texttt{rc-charging}.\par}

\end{document}
